% from ega2-1-fr.tex
\oldpage[II]{200}
Supposons d'abord $Y=\operatorname{Spec}(A)$ affine,
$\mathcal{S}=\widetilde{S}$, où $S$ est une $A$-algèbre graduée à degrés
positifs~; alors $\Gamma'_*(\mathcal{F}')=\widetilde{M}$, où
$M=\bigoplus_{n\in\mathbf{Z}}\Gamma(X',\mathcal{F}'_n)$ est un $S$-module
gradué. Soit $f\in S_d$ tel que $D_+(f)\subset X'$~; par définition
(\hyperref[II.2.6.2-fr]{2.6.2}), $\alpha_d(f)$ restreinte à $D_+(f)$ est la
section de $\mathcal{O}_X(d)$ au-dessus de $D_+(f)$ correspondant à
l'élément $f/1$ de $(S(d))_{(f)}$ et est par suite inversible~; il en est donc
de même de $\alpha_d(f^n)$ pour tout $n>0$. On en conclut aussitôt que
l'on définit un $S_f$-homomorphisme (de degré $0$) de modules gradués
$\beta_f:M_f\to\Gamma(D_+(f),\mathcal{F}')$ en faisant correspondre à tout
élément $z/f^n\in M_f$ (où $z\in M$), la section
$(z|D_+(f))(\alpha_d(f^n)|D_+(f))^{-1}$ de $\mathcal{F}'$ au-dessus de
$D_+(f)$. En outre, on a un diagramme commutatif correspondant à
(\hyperref[II.2.6.4.1-fr]{2.6.4.1}), d'où la définition de $\beta'$ dans ce
cas. Pour passer au cas général, il faut considérer une $A$-algèbre $A'$,
la $A'$-algèbre graduée $S'=S\otimes_A A'$, et utiliser un diagramme
commutatif analogue à (\hyperref[II.2.8.13.2-fr]{2.8.13.2})~; nous laissons
les détails aux lecteurs.
\end{env}

\begin{proposition}[8.14.8]
\label{II.8.14.8-fr}
Si $X'$ est un ouvert de $X=\operatorname{Proj}(\mathcal{S})$ tel que
$q':X'\to Y$ soit quasi-compact, l'homomorphisme $\beta'$ défini dans
(\hyperref[II.8.14.7-fr]{8.14.7}) est bijectif.
\end{proposition}

\begin{proof}
On peut évidemment se restreindre au cas où $Y$ est affine et tout revient à
prouver (avec les notations de (\hyperref[II.8.14.7-fr]{8.14.7})) que
l'homomorphisme $\beta_f:M_f\to\Gamma(D_+(f),\mathcal{F}')$ est un
isomorphisme. Or, lorsqu'on remplace $f$ par une de ses puissances, on ne
change pas $D_+(f)$ ni $\beta_f$~; comme $X'$ est \emph{quasi-compact} par
hypothèse, on peut donc toujours supposer, en vertu de
(\hyperref[II.8.14.4-fr]{8.14.4}), que le faisceau $\mathcal{O}_X(d)$ est
\emph{inversible}. Comme $X'$ est un schéma (puisque $q'$ est séparé), la
proposition n'est autre alors que (\hyperref[I.9.3.1-fr]{I, 9.3.1}).
\end{proof}

\begin{corollary}[8.14.9]
\label{II.8.14.9-fr}
Sous les hypothèses de (\hyperref[II.8.14.8-fr]{8.14.8}), tout
$\mathcal{S}_{X'}$-Module gradué quasi-cohérent $\mathcal{F}'$ est isomorphe
à un $\mathcal{S}_{X'}$-Module gradué de la forme
$\operatorname{Proj}'(\mathcal{M})$, où $\mathcal{M}$ est un
$\mathcal{S}$-Module gradué quasi-cohérent. Si en outre $\mathcal{F}'$ est de
type fini, et si on suppose que $Y$ est un schéma quasi-compact ou un
préschéma dont l'espace sous-jacent est noethérien, on peut supposer
$\mathcal{M}$ de type fini.
\end{corollary}

\begin{proof}
La démonstration à partir de (\hyperref[II.8.14.8-fr]{8.14.8}) suit
exactement la marche de celle de (\hyperref[II.3.4.5-fr]{3.4.5}) à partir de
(\hyperref[II.3.4.4-fr]{3.4.4}), et nous en laissons les détails au lecteur.
\end{proof}

\begin{proposition}[8.14.10]
\label{II.8.14.10-fr}
Sous les hypothèses de (\hyperref[II.8.14.7-fr]{8.14.7}), soient
$\mathcal{M}$ un $\mathcal{S}$-Module gradué quasi-cohérent,
$\mathcal{F}'$ un $\mathcal{S}_{X'}$-Module gradué quasi-cohérent~; les
homomorphismes composés
\begin{equation}
\label{II.8.14.10.1-fr}
  \operatorname{Proj}'(\mathcal{M})
  \xrightarrow{\operatorname{Proj}'(\alpha')}
  \operatorname{Proj}'(\Gamma'_*(\operatorname{Proj}'(\mathcal{M})))
  \xrightarrow{\beta'}\operatorname{Proj}'(\mathcal{M})
\tag{8.14.10.1}
\end{equation}
\begin{equation}
\label{II.8.14.10.2-fr}
  \Gamma'_*(\mathcal{F}')
  \xrightarrow{\alpha'}
  \Gamma'_*(\operatorname{Proj}'(\Gamma'_*(\mathcal{F}')))
  \xrightarrow{\Gamma'_*(\beta')}\Gamma'_*(\mathcal{F}')
\tag{8.14.10.2}
\end{equation}
sont les isomorphismes identiques.
\end{proposition}

\begin{proof}
La question est locale sur $Y$ et la vérification se fait comme dans
(\hyperref[II.2.6.5-fr]{2.6.5})~; nous en laissons les détails au lecteur.
\end{proof}

\begin{remark}[8.14.11]
\label{II.8.14.11-fr}
Au chapitre III (\hyperref[III.2.3.1-fr]{III, 2.3.1}), nous verrons que
lorsque $Y$ est \emph{localement noethérien} et $\mathcal{S}$ une
$\mathcal{O}_Y$-Algèbre graduée quasi-cohérente \emph{de type fini} (auquel
cas

